\documentclass[12pt,a4paper]{article}
% Drake_preamble.tex - LaTeX Preamble Template
% 作者: 謝慕揚, MD, PhD, FESC
% 
% 使用說明:
% 1. 表格請使用 longtable，不要有垂直分隔線 |
% 2. 表格內換行使用 \newline，不要使用 \\
% 3. 藥物名稱、檢驗、檢查請維持英文
% 4. Reference 格式請使用 \href 加上驗證過的 DOI

% Including essential packages for Chinese typesetting
\usepackage{xeCJK}
\usepackage{fontspec}
% Configuring Chinese font with Noto Serif CJK TC, fallback to SimSun
\IfFontExistsTF{Noto Serif CJK TC}{
  \setCJKmainfont{Noto Serif CJK TC}
  \setCJKsansfont{Noto Serif CJK TC}
  \setCJKmonofont{Noto Serif CJK TC}
}{\setCJKmainfont{SimSun}
  \setCJKsansfont{SimSun}
  \setCJKmonofont{SimSun}
}

% Including other necessary packages
\usepackage{geometry}
\usepackage{titlesec}
\usepackage{enumitem}
\usepackage{xcolor}
\usepackage{colortbl}
\usepackage{tcolorbox}
\usepackage{booktabs}
\usepackage{longtable}
\usepackage{array}
\usepackage{graphicx}
\usepackage{tikz}
\usepackage{fancyhdr}
\usepackage{multicol}
\usepackage{datetime2}
\usepackage{hyperref}
\usepackage{mdframed}
\usepackage{amsmath}
\usepackage{amssymb}
\usepackage{multirow}

\hypersetup{
    colorlinks=true,
    linkcolor=emergency_blue,
    citecolor=emergency_blue,
    urlcolor=emergency_blue,
    pdfborder={0 0 0}
}

% Defining page geometry
\geometry{
  top=2.5cm,
  bottom=2.5cm,
  left=2cm,
  right=2cm,
  headheight=25pt
}

% Adjusting topmargin to compensate for increased headheight
\addtolength{\topmargin}{-10pt}

% Defining colors
\definecolor{emergency_red}{RGB}{186,24,27}
\definecolor{emergency_blue}{RGB}{0,114,188}
\definecolor{emergency_gray}{RGB}{120,120,120}
\definecolor{highlight_yellow}{RGB}{255,250,205}

% Setting title formats
\titleformat{\section}[block]{\Large\bfseries\color{emergency_red}}{\thesection}{10pt}{}
\titleformat{\subsection}[block]{\large\bfseries\color{emergency_blue}}{\thesubsection}{8pt}{}
\titleformat{\subsubsection}[block]{\normalsize\bfseries\color{emergency_gray}}{\thesubsubsection}{6pt}{}

% Defining custom environment for important notes
\newmdenv[
    linecolor=emergency_red,
    backgroundcolor=highlight_yellow,
    linewidth=2pt,
    topline=true,
    bottomline=true,
    leftline=true,
    rightline=true,
    innertopmargin=10pt,
    innerbottommargin=10pt,
    innerleftmargin=10pt,
    innerrightmargin=10pt
]{importantbox}

% Defining note command with tcolorbox
\newcommand{\note}[1]{
    \par
    \noindent
    \begin{tcolorbox}[
        colback=emergency_blue!5!white,
        colframe=emergency_blue,
        title=\textbf{重點提醒},
        fonttitle=\bfseries,
        boxrule=0.5pt,
        arc=3pt,
        left=6pt,
        right=6pt,
        top=6pt,
        bottom=6pt
    ]
    #1
    \end{tcolorbox}
}

% Key findings box
\newcommand{\keyfinding}[1]{
    \par
    \noindent
    \begin{tcolorbox}[
        colback=yellow!10!white,
        colframe=orange!75!black,
        title=\textbf{核心發現摘要},
        fonttitle=\bfseries,
        boxrule=0.5pt,
        arc=3pt
    ]
    #1
    \end{tcolorbox}
}

% Defining custom date format for ISO 8601
\DTMnewdatestyle{iso}{
  \renewcommand{\DTMdisplaydate}[4]{\number##1-\number##2-\number##3}
  \renewcommand{\DTMDisplaydate}[4]{\number##1-\number##2-\number##3}
}
\DTMsetdatestyle{iso}
\newcommand{\mydate}{\DTMtoday}


% Custom header for this document
\pagestyle{fancy}
\fancyhf{}
\fancyhead[L]{\textcolor{emergency_red}{Intensive Care Medicine 教學講義}}
\fancyhead[R]{\textcolor{emergency_blue}{Clinical Governance in ICU}}
\fancyfoot[C]{\textcolor{emergency_gray}{\thepage}}
\renewcommand{\headrulewidth}{1pt}
\renewcommand{\headrule}{\hbox to\headwidth{\color{emergency_red}\leaders\hrule height \headrulewidth\hfill}}

\begin{document}

\begin{center}
{\LARGE\bfseries\textcolor{emergency_red}{Intensive Care Medicine Editorial}}\\[0.5cm]
{\Large\textbf{Clinical Governance in Intensive Care Medicine}}\\[0.3cm]
{\large 加護病房的臨床治理}\\[0.5cm]
{\normalsize Intensive Care Med 2024;50:2154-2157. DOI: 10.1007/s00134-024-07653-8}\\[0.3cm]
{\normalsize 整理：謝慕揚 MD, PhD, FESC \quad 日期：\mydate}
\end{center}

\keyfinding{
\textbf{核心訊息：}Clinical governance (CG) 是醫療機構（包括 ICU）持續改善服務並提供安全、有效、具同理心照護的責任機制。

\textbf{實用重點：}
\begin{itemize}
    \item \textbf{五大文化基礎：}Systems awareness、Teamwork、Communication、Ownership、Leadership
    \item \textbf{關鍵工具：}Checklists、Auditing、Incident reporting、M\&M conferences、Simulation
    \item \textbf{整合效果：}結合多種 CG 元素可產生 synergistic effect，比單一元素更有效改善病人預後
\end{itemize}
}

\tableofcontents
\newpage

%==============================================================
\section{文獻概述}
%==============================================================

本篇為 Intensive Care Medicine 的 Editorial，由義大利 Humanitas Research Hospital 的 Carenzo 等人撰寫，介紹 clinical governance 的概念及其在重症醫學的實務應用。

\subsection{作者資訊}
\textbf{通訊作者：}Luca Carenzo, MD\\
\textbf{機構：}Department of Anesthesia and Intensive Care Medicine, IRCCS Humanitas Research Hospital, Rozzano, Italy

%==============================================================
\section{Clinical Governance 概述}
%==============================================================

\subsection{定義}

\begin{itemize}
    \item \textbf{Clinical Governance (CG)：}醫療機構（包括 ICU）對持續改善服務並提供 safe, effective and compassionate care 所負的責任
    \item Patient safety 是 healthcare improvement 的首要焦點
    \item ICU 的複雜性增加 medical errors 風險
\end{itemize}

\subsection{Temple Model}

\note{
\textbf{Clinical Governance 的神殿模型：}

\textbf{基礎（五大文化元素）：}
\begin{enumerate}
    \item Systems awareness（系統意識）
    \item Teamwork（團隊合作）
    \item Communication（溝通）
    \item Ownership（責任歸屬）
    \item Leadership（領導力）
\end{enumerate}

\textbf{支柱：}Clinical effectiveness、Clinical audit、Risk management、Patient and public involvement、Staff management、Education and training、Openness
}

%==============================================================
\section{CG 領域與實用工具}
%==============================================================

\begin{longtable}{p{4cm}p{10cm}}
\toprule
\textbf{領域} & \textbf{工具} \\
\midrule
\endfirsthead

\toprule
\textbf{領域} & \textbf{工具} \\
\midrule
\endhead

\bottomrule
\endfoot

\textbf{Education, Training \& Staff} &
• Multidisciplinary simulation\newline
• Staff well-being programs\newline
• Enhancing teamwork and communication \\
\midrule
\textbf{Effective Care} &
• Benchmarking\newline
• Checklists\newline
• Auditing \\
\midrule
\textbf{Risk Management \& Openness} &
• Morbidity and mortality meetings\newline
• Incident reporting\newline
• Clinical governance meetings \\
\midrule
\textbf{Patient and Family Centered Care} &
• Family involvement in patient care\newline
• Unrestricted visitation\newline
• Family and patient experience surveys \\
\bottomrule
\end{longtable}

%==============================================================
\section{Patient and Family Centered Care (PFCC)}
%==============================================================

\subsection{定義}

\begin{itemize}
    \item 以病人及家屬的 needs, values, preferences 為中心的照護模式
    \item 主動將家屬納入 decision-making 及 care process
    \item 特別是當病人無法自行溝通時
\end{itemize}

\subsection{有效介入措施}

\begin{longtable}{p{4cm}p{10cm}}
\toprule
\textbf{介入措施} & \textbf{說明} \\
\midrule
\endfirsthead

\toprule
\textbf{介入措施} & \textbf{說明} \\
\midrule
\endhead

\bottomrule
\endfoot

\textbf{Unrestricted Visitation} & 鼓勵家屬在場，不限制探視時間 \\
\midrule
\textbf{Daily Care Engagement} & 讓家屬參與日常照護及 medical rounds \\
\midrule
\textbf{Psychological Support} & ICU diaries、communication enhancements \\
\bottomrule
\end{longtable}

\subsection{PFCC 的效益}

\begin{itemize}
    \item 降低 anxiety
    \item 提升 patient and family satisfaction
    \item 降低 delirium 風險
    \item 降低 post-traumatic stress disorder 風險
    \item 減少 ICU staff burnout
\end{itemize}

%==============================================================
\section{Effective Teamwork and Communication}
%==============================================================

\begin{longtable}{p{4cm}p{10cm}}
\toprule
\textbf{工具} & \textbf{說明} \\
\midrule
\endfirsthead

\toprule
\textbf{工具} & \textbf{說明} \\
\midrule
\endhead

\bottomrule
\endfoot

\textbf{Interprofessional Ward Rounds} &
• 各專業人員共同參與\newline
• 做出全面性決策\newline
• 增進 patient safety、減少 errors \\
\midrule
\textbf{Structured Handovers} &
• 傳遞資訊及責任\newline
• 使用 memory aids 及 standardized protocols\newline
• 最小化 errors \\
\midrule
\textbf{Safety Huddles} &
• 班前簡短的跨專業會議\newline
• 識別潛在風險\newline
• 確保所有成員掌握 critical patient safety information \\
\bottomrule
\end{longtable}

%==============================================================
\section{Staff Management and Well-Being}
%==============================================================

\subsection{重要性}

\begin{itemize}
    \item 確保團隊有效招募、訓練及支持
    \item 營造正向工作環境
    \item 處理 staff development needs
    \item 因應 staff shortages 及 burnout 挑戰
\end{itemize}

\subsection{Iatroref Study 發現}

\note{
\textbf{重要發現：}ICU staff 的 depression symptoms 與 medical errors 風險增加顯著相關。

\textbf{臨床意涵：}有效的 CG interventions 應促進 staff well-being 並降低 burnout 及 depression rates。
}

%==============================================================
\section{Checklists}
%==============================================================

\subsection{功能與效益}

\begin{itemize}
    \item 作為 cognitive aids，確保遵循 best practices
    \item 標準化流程、改善團隊表現、減少 ICU errors
    \item 類似航空業的 checklist 應用
\end{itemize}

\subsection{常見 ICU Checklists}

\begin{longtable}{p{5cm}p{9cm}}
\toprule
\textbf{類別} & \textbf{內容} \\
\midrule
\endfirsthead

\toprule
\textbf{類別} & \textbf{內容} \\
\midrule
\endhead

\bottomrule
\endfoot

\textbf{有正面效果的 Checklists} &
• Delirium screening\newline
• CLABSI prevention\newline
• Airway management \\
\midrule
\textbf{Intensive Care Society Checklists} &
• Central venous catheter insertion\newline
• Nasogastric tube insertion\newline
• Chest drains\newline
• Bronchoscopies\newline
• Intubations\newline
• Tracheostomies \\
\bottomrule
\end{longtable}

%==============================================================
\section{Auditing and Benchmarking}
%==============================================================

\subsection{Auditing 定義}

\begin{itemize}
    \item 系統性審查臨床實務，確保符合既定標準
    \item 應視為成長機會，而非評判性評估
    \item 提供 feedback 以強化 guideline compliance 及促進改變
\end{itemize}

\subsection{ICU 主要 Audit 領域}

\begin{itemize}
    \item Ventilation
    \item Ventilator-associated pneumonia
    \item Intubation
    \item Tracheostomy care
    \item Venous access bundles compliance
    \item Renal replacement therapy
\end{itemize}

\subsection{Auditing 循環}

\begin{enumerate}
    \item 收集資料
    \item 評估資料
    \item 重新收集資料以測量改善
    \item 在 CG meetings 分享結果
\end{enumerate}

\note{
\textbf{Cochrane Review 發現：}Audit and feedback 可帶來 small but significant improvements in practice。
}

\subsection{Benchmarking}

\begin{itemize}
    \item 使用標準化指標比較 ICU 表現
    \item 可進行內部及跨機構比較
\end{itemize}

\textbf{常用指標：}
\begin{itemize}
    \item Mortality rates
    \item Lengths of stay
    \item Readmission rates
    \item Adherence to care protocols
    \item Patient satisfaction
\end{itemize}

%==============================================================
\section{Incident Reporting}
%==============================================================

\subsection{定義}

公開揭露並記錄臨床實務中發生的 adverse events、errors 及 near misses。

\subsection{Incident Reporting Systems (IRS) 的功能}

\begin{itemize}
    \item 收集事件
    \item 提供資料以反思事件發生原因
    \item 思考預防方法
    \item 為 training programs 及 policy changes 提供資訊
\end{itemize}

\subsection{成功實施的關鍵}

\begin{itemize}
    \item 使用 web-based 或 paper-based 平台
    \item 確保 confidentiality，無報復風險
    \item 教育 staff 了解 reporting system 的重要性
    \item 減少因 fear of blame 造成的 underreporting
    \item 提供及時 feedback
\end{itemize}

%==============================================================
\section{Morbidity and Mortality (M\&M) Conferences}
%==============================================================

\subsection{功能}

\begin{itemize}
    \item 系統性審查 adverse events 及 patient deaths
    \item 識別 preventable errors 或 systemic issues
    \item 增進 patient safety 及 care quality
\end{itemize}

\subsection{會議特點}

\begin{itemize}
    \item 鼓勵 staff participation
    \item 連結 case presentations、discussions 及 action items
    \item 識別 systemic improvements 機會
\end{itemize}

\subsection{M\&M Meetings 的目標}

\begin{itemize}
    \item Quality improvement
    \item Education
    \item Legal risk management
    \item Peer support
\end{itemize}

\note{
\textbf{研究發現：}Caregivers 對 adverse events preventability 的評估與 ICU organizational culture 之間有強烈關聯。

\textbf{結論：}Error communication 引發的教育及改變是改善 ICU 表現的第一步。
}

%==============================================================
\section{Multidisciplinary Simulation}
%==============================================================

\subsection{定義}

\textit{``Simulation is a technique—not a technology—to replace or amplify real experiences with guided experiences that evoke or replicate substantial aspects of the real world in a fully interactive manner.''} — David Gaba

\subsection{In-Situ Simulation}

\begin{longtable}{p{4cm}p{10cm}}
\toprule
\textbf{項目} & \textbf{說明} \\
\midrule
\endfirsthead

\toprule
\textbf{項目} & \textbf{說明} \\
\midrule
\endhead

\bottomrule
\endfoot

\textbf{定義} & 在 ICU 內使用真實設備進行的模擬訓練 \\
\midrule
\textbf{優點} &
• 訓練與真實情境相關\newline
• 增進 technical 及 non-technical skills\newline
• 發現 local system errors 及 latent safety threats\newline
• 發現 faulty equipment 或 inadequate resources \\
\midrule
\textbf{情境來源} &
• 醫院 incident reporting system\newline
• 真實 ICU cases \\
\midrule
\textbf{常見情境} &
• Airway failure\newline
• Accidental extubation\newline
• Medication errors\newline
• Sepsis recognition \\
\bottomrule
\end{longtable}

\subsection{Debriefing 的重要性}

\begin{itemize}
    \item 將 simulation experiences 轉化為 learning opportunities
    \item 透過 structured reflection 及 feedback 進行
    \item 結果應記錄並分享給相關 stakeholders
\end{itemize}

%==============================================================
\section{Clinical Governance Meetings}
%==============================================================

\subsection{功能}

\begin{itemize}
    \item Information sharing 的核心場合
    \item 定期舉行的跨專業會議
\end{itemize}

\subsection{討論內容}

\begin{itemize}
    \item ICU guidelines
    \item Innovations
    \item Complaints
    \item Incident reports
    \item Audits
    \item Morbidity and mortality cases
\end{itemize}

\subsection{參與人員}

\begin{itemize}
    \item 所有 ICU professionals（開放參與）
    \item Hospital management
    \item External stakeholders
    \item Patient representatives
\end{itemize}

\note{
\textbf{重要原則：}Meeting results 應予以記錄，作為 process 的 evidence。
}

%==============================================================
\section{Take-Home Message}
%==============================================================

\begin{enumerate}
    \item \textcolor{red}{\textbf{系統性應用：}}系統性應用 CG 工具可協助 ICU 識別風險並促進 transparency、accountability 及 collective learning 的文化
    \item \textcolor{red}{\textbf{整合效果：}}結合多種 CG 元素（如 checklists、staff education、regular feedback）可產生 synergistic effect
    \item \textcolor{red}{\textbf{Staff Well-Being：}}Staff depression 與 medical errors 相關，應積極促進 staff well-being
    \item \textcolor{red}{\textbf{持續改善：}}透過 auditing、incident reporting、M\&M conferences 及 simulation 持續改善照護品質
\end{enumerate}

%==============================================================
\section{縮寫對照表}
%==============================================================

\begin{longtable}{p{3cm}p{11cm}}
\toprule
\textbf{縮寫} & \textbf{全名} \\
\midrule
\endfirsthead

\toprule
\textbf{縮寫} & \textbf{全名} \\
\midrule
\endhead

\bottomrule
\endfoot

CG & Clinical Governance (臨床治理) \\
ICU & Intensive Care Unit (加護病房) \\
PFCC & Patient and Family Centered Care \\
IRS & Incident Reporting Systems \\
M\&M & Morbidity and Mortality \\
CLABSI & Central Line-Associated Bloodstream Infection \\
\bottomrule
\end{longtable}

%==============================================================
\section{參考文獻}
%==============================================================

\begin{enumerate}
    \item Carenzo L, Costantini E, Cecconi M. Clinical governance in intensive care medicine. \href{https://doi.org/10.1007/s00134-024-07653-8}{\textit{Intensive Care Med}. 2024;50:2154-2157.}

    \item Pronovost P, Needham D, Berenholtz S, et al. An intervention to decrease catheter-related bloodstream infections in the ICU. \href{https://doi.org/10.1056/NEJMoa061115}{\textit{N Engl J Med}. 2006;355:2725-2732.}

    \item Garrouste-Orgeas M, Perrin M, Soufir L, et al. The Iatroref study: medical errors are associated with symptoms of depression in ICU staff but not burnout or safety culture. \href{https://doi.org/10.1007/s00134-014-3601-4}{\textit{Intensive Care Med}. 2015;41:273-284.}

    \item Rhodes A, Moreno RP, Azoulay E, et al. Prospectively defined indicators to improve the safety and quality of care for critically ill patients: a report from the Task Force on Safety and Quality of the European Society of Intensive Care Medicine (ESICM). \href{https://doi.org/10.1007/s00134-011-2462-3}{\textit{Intensive Care Med}. 2012;38:598-605.}

    \item Ivers N, Jamtvedt G, Flottorp S, et al. Audit and feedback: effects on professional practice and healthcare outcomes. \href{https://doi.org/10.1002/14651858.CD000259.pub3}{\textit{Cochrane Database Syst Rev}. 2012;6:CD000259.}
\end{enumerate}

\end{document}
